% Generated by scripts/build_dashboard.py
\begin{itemize}
\item 2000: Analysis of the genome sequence of the flowering plant Arabidopsis thaliana
\item 2000: An RNA-dependent RNA polymerase gene in Arabidopsis is required for posttranscriptional gene silencing mediated by a transgene but not by a virus
\item 2000: Sequence and analysis of chromosome 3 of the plant Arabidopsis thaliana
\item 2000: Sequence and analysis of chromosome 5 of the plant Arabidopsis thaliana
\item 2000: Arabidopsis genome analysis as exemplified by analysis of chromosome 4
\item 2002: Full-length messenger RNA sequences greatly improve genome annotation
\item 2002: MIPS: a database for genomes and protein sequences
\item 2002: MIPS Arabidopsis thaliana Database (MAtDB): an integrated biological knowledge resource based on the first complete plant genome
\item 2003: Expressed sequence tag analysis in Cycas, the most primitive living seed plant
\item 2003: The genome sequence of the filamentous fungus Neurospora crassa
\item 2003: Snipping polymorphisms from large EST collections in barley (Hordeum vulgare L.)
\item 2003: Expressed sequence tags: alternative or complement to whole genome sequences?
\item 2003: Sputnik: a database platform for comparative plant genomics
\item 2004: Characterization of the maize endosperm transcriptome and its comparison to the rice genome
\item 2004: Genome-wide in silico mapping of scaffold/matrix attachment regions in Arabidopsis suggests correlation of intragenic scaffold/matrix attachment regions with gene expression
\item 2004: Large-scale analysis of the barley transcriptome based on expressed sequence tags
\item 2005: Identification of NdhL and Ssl1690 (NdhO) in NDH-1L and NDH-1M complexes of Synechocystis sp. PCC 6803
\item 2005: EST analysis in Ginkgo biloba: an assessment of conserved developmental regulators and gymnosperm-specific genes
\item 2005: Gene expression and metabolite profiling of Populus euphratica growing in the Negev desert
\item 2005: Support vector machines for separation of mixed plant-pathogen EST collections based on codon usage
\item 2005: A 6374 unigene set corresponding to low abundance transcripts expressed following fertilization in Solanum chacoense and characterization of 30 receptor-like kinases
\item 2005: Analysis of the floral transcriptome uncovers new regulators of organ determination and gene families related to flower organ differentiation in Gerbera hybrida (Asteraceae)
\item 2005: Eclair---a web service for unravelling species origin of sequences sampled from mixed host interfaces
\item 2005: openSputnik---a database to ESTablish comparative plant genomics using unsaturated sequence collections
\item 2005: PlantMarkers---a database of predicted molecular markers from plants
\item 2005: Maintenance of ancestral complexity and non-metazoan genes in two basal cnidarians
\item 2006: A wing expressed sequence tag resource for Bicyclus anynana butterflies, an evo-devo model
\item 2006: Genetic screen for signal peptides in Hydra reveals novel secreted proteins and evidence for non-classical protein secretion
\item 2006: Prospecting for pig single nucleotide polymorphisms in the human genome: have we struck gold?
\item 2006: Construction, database integration, and application of an Oenothera EST library
\item 2006: Spatiotemporal expression control correlates with intragenic scaffold matrix attachment regions (S/MARs) in Arabidopsis thaliana
\item 2007: Separation of sequences from host-pathogen interface using triplet nucleotide frequencies
\item 2007: Expression profiling of the lignin biosynthetic pathway in Norway spruce using EST sequencing and real-time RT-PCR
\item 2007: Characterization of an 18,166 EST dataset for cassava (Manihot esculenta Crantz) enriched for drought-responsive genes
\item 2007: Proteomic screen in the simple metazoan Hydra identifies 14-3-3 binding proteins implicated in cellular metabolism, cytoskeletal organisation and Ca2+ signalling
\item 2007: Analysis of EST sequences suggests recent origin of allotetraploid colonial and creeping bentgrasses
\item 2008: Microarray analysis identifies candidate genes for key roles in coral development
\item 2009: Comparative mapping of DNA sequences in rye (Secale cereale L.) in relation to the rice genome
\item 2009: Optimization and comparison of different methods for RNA isolation for cDNA library construction from the reindeer lichen Cladonia rangiferina
\item 2012: Draft genome sequences of Helicobacter pylori isolates from Malaysia, cultured from patients with functional dyspepsia and gastric cancer
\item 2012: Characterization of a transcriptome from a non-model organism, Cladonia rangiferina, the grey reindeer lichen, using high-throughput next generation sequencing and EST sequence data
\item 2013: Whole transcriptome characterization of the effects of dehydration and rehydration on Cladonia rangiferina, the grey reindeer lichen
\item 2014: Programmed ribosomal frameshift alters expression of West Nile virus genes and facilitates virus replication in birds and mosquitoes
\item 2016: Gender-specific molecular and clinical features underlie malignant pleural mesothelioma
\item 2016: Human mesenchymal stem cells alter the gene profile of monocytes from patients with type 2 diabetes and end-stage renal disease
\item 2017: Bioinformatics in the plant genomic and phenomic domain: the German contribution to resources, services and perspectives
\item 2018: Rapid analysis of metagenomic data using signature-based clustering
\item 2018: Gene networks associated with non-syndromic intellectual disability
\end{itemize}
